%! Characteristics of Radical Functions — Unit 6, Lesson 6.0 — Homework (Answer Key)
\documentclass[10pt]{article}
\usepackage{algebra2-article}
\usepackage{algebra2-key}   % pulls in -boxes; do NOT also load -boxes
\newcommand{\TallMath}[1]{$\displaystyle #1\rule[-1.4em]{0pt}{3.2em}$}
\newcommand{\lgrid}[4]{%
  \draw[step=1, linegray!40, very thin] (#1,#3) grid (#2,#4);
  \draw[->, semithick, charcoal] (#1,0)--(#2,0) node[below right, font=\scriptsize]{$x$};
  \draw[->, semithick, charcoal] (0,#3)--(0,#4) node[left, font=\scriptsize]{$y$};}
\newcommand{\cbrtpos}[2]{\draw[very thick, forest, domain=#1:#2, samples=70, smooth]
  plot (\x,{pow(\x,1/3)+1});}
\newcommand{\cbrtneg}[2]{\draw[very thick, forest, domain=#1:#2, samples=70, smooth]
  plot (\x,{-pow(-1*(\x),1/3)+1});}

\begin{document}

\pageheader{Unit 6, Lesson 6.0}{Homework --- Answer Key}
\namedateperiod

\vspace{0.10in}

\begin{notesbox}{Practice}
Show how you read each feature. Write every domain and range in \textbf{interval notation}.

\begin{enumerate}[leftmargin=1.9em, itemsep=12pt]
  \item \textbf{From the equation only.} Use the index and the radicand --- no graphing.
    \begin{center}
    \renewcommand{\arraystretch}{1.6}
    \begin{tabularx}{\linewidth}{@{}l X X X X@{}}
    \hline
    \textbf{Function} & \textbf{Index: even/odd} & \textbf{Condition on the radicand} & \textbf{Domain} & \textbf{Opens left, right, or both ways?} \\
    \hline
    $\sqrt{x+6}$ & \ans{even} & \ans{$x+6\ge 0$} & \ans{$[-6,\infty)$} & \ans{right} \\
    $\sqrt[3]{x-4}$ & \ans{odd} & \ans{none} & \ans{$(-\infty,\infty)$} & \ans{both} \\
    $\sqrt{3x-9}$ & \ans{even} & \ans{$3x-9\ge 0$} & \ans{$[3,\infty)$} & \ans{right} \\
    $\sqrt{8-4x}$ & \ans{even} & \ans{$8-4x\ge 0$} & \ans{$(-\infty,2]$} & \ans{left} \\
    \hline
    \end{tabularx}
    \end{center}

  \item \textbf{A square root graph that falls.} The graph is $m(x) = -\sqrt{x}+3$.
    \begin{center}
    \begin{tikzpicture}[scale=0.44]
      \lgrid{-2}{11}{-2}{4}
      \begin{scope}
        \clip (-2,-2) rectangle (11,4);
        \draw[very thick, forest, domain=0:11, samples=100, smooth] plot (\x,{3-sqrt(\x)});
      \end{scope}
      \fill[forest] (0,3) circle (4pt) node[above left, font=\scriptsize]{$(0,3)$};
      \fill[charcoal] (1,2) circle (3pt);
      \fill[charcoal] (4,1) circle (3pt) node[above right, font=\scriptsize]{$(4,1)$};
      \fill[charcoal] (9,0) circle (3pt) node[below right, font=\scriptsize]{$(9,0)$};
    \end{tikzpicture}
    \end{center}
    Domain: \ans{$[0,\infty)$} \quad Range: \ans{$(-\infty,3]$} \quad Endpoint: \ans{$(0,3)$} \\[3pt]
    $x$-intercept: \ans{$(9,0)$} \quad $y$-intercept: \ans{$(0,3)$} \\[3pt]
    Increasing or decreasing on its domain? \ans{decreasing} \\[3pt]
    Absolute \ans{maximum} of \ans{3} at $x=$ \ans{0}; \; no absolute \ans{minimum} \\[3pt]
    End behavior: as $x\to+\infty$, $m(x)\to$ \ans{$-\infty$}. \; As $x\to-\infty$:
    \ans{the graph stops at $x=0$}

  \item \textbf{A cube root graph.} The graph is $n(x) = \sqrt[3]{x}+1$.
    \begin{center}
    \begin{tikzpicture}[scale=0.44]
      \lgrid{-9}{9}{-2}{4}
      \begin{scope}
        \clip (-9,-2) rectangle (9,4);
        \cbrtpos{0}{9}
        \cbrtneg{-9}{0}
      \end{scope}
      \fill[charcoal] (-8,-1) circle (3pt) node[below right, font=\scriptsize]{$(-8,-1)$};
      \fill[charcoal] (-1,0) circle (3pt);
      \fill[charcoal] (0,1) circle (3pt) node[above left, font=\scriptsize]{$(0,1)$};
      \fill[charcoal] (8,3) circle (3pt) node[below right, font=\scriptsize]{$(8,3)$};
    \end{tikzpicture}
    \end{center}
    Domain: \ans{$(-\infty,\infty)$} \quad Range: \ans{$(-\infty,\infty)$} \quad Endpoint:
    \ans{none} \\[3pt]
    $x$-intercept: \ans{$(-1,0)$} \quad $y$-intercept: \ans{$(0,1)$} \\[3pt]
    Increasing or decreasing? \ans{increasing} \quad Absolute extrema: \ans{none} \\[3pt]
    End behavior: as $x\to+\infty$, $n(x)\to$ \ans{$+\infty$}; \; as $x\to-\infty$,
    $n(x)\to$ \ans{$-\infty$} \\[3pt]
    The \emph{parent} $y=\sqrt[3]{x}$ is symmetric about the origin. Is $n$? \ans{No} \;
    Why or why not? \ans{The whole graph was shifted up $1$, so its center of symmetry moved from
    $(0,0)$ to $(0,1)$. Check: $n(-8)=-1$, but $-n(8)=-3$.}

  \item \textbf{Explain.} A square root graph always has exactly one absolute extremum, but a cube
        root graph has none at all. Explain why, using the word \textbf{endpoint}.
        \ansline{A square root graph has an \textbf{endpoint} and leaves it in only one direction, so
        that single point is either the highest or the lowest point on the entire graph --- an
        absolute maximum if the arm falls away from it, an absolute minimum if the arm climbs.\\[3pt]
        A cube root graph has no endpoint. It runs to $-\infty$ on the left and $+\infty$ on the
        right, so no point on it is highest or lowest, and it has neither kind of extremum.}
\end{enumerate}
\end{notesbox}

\begin{extensionbox}
\textbf{Reading a model in context.} Accident investigators estimate how fast a car was travelling
from the length of its skid marks. On dry asphalt, a skid of $d$ feet means the car was going about
\[
  s(d) = \sqrt{24d} \ \text{ miles per hour.}
\]
\begin{center}
\renewcommand{\arraystretch}{1.25}
\begin{tabular}{|c|c|c|c|c|c|}
\hline
$d$ (feet of skid) & $6$ & $24$ & $54$ & $96$ & $150$ \\ \hline
$s(d)$ (mph) & $12$ & $24$ & $36$ & $48$ & $60$ \\ \hline
\end{tabular}
\end{center}
\begin{enumerate}[label=(\alph*), leftmargin=1.8em, itemsep=5pt]
  \item Find the domain of $s$ from the radicand. Does the context restrict it any further?
        \ans{Even index, so $24d\ge 0$, giving $d\ge 0$ --- the domain is $[0,\infty)$. The context agrees exactly: a skid mark cannot have negative length.}
  \item The graph of $s$ has an endpoint at $(0,0)$. What does that point mean about a car that
        leaves no skid marks at all?
        \ans{The formula estimates $0$ mph --- no skid, no speed. That is also the model's honest limit: a car with anti-lock brakes can be moving fast and still leave no skid marks, so $s$ only applies when there \emph{are} marks to measure.}
  \item A $24$-foot skid means about $24$ mph. To double that estimate to $48$ mph, the table says the
        skid must be $96$ feet --- \emph{four} times as long, not twice. Explain why, using the shape
        of a square root graph.
        \ans{A square root graph flattens as it climbs, so each extra mph costs more and more skid. Since square-rooting undoes squaring, doubling the output requires multiplying the input by $2^2=4$.}
  \item Solve $s=\sqrt{24d}$ for $d$. Which function family did you get, and what is its relationship
        to $s(d)$?
        \ans{Square both sides: $s^2=24d$, so $d=\frac{s^2}{24}$ --- a \textbf{quadratic}. It is the \textbf{inverse} of $s(d)$: one turns skid length into speed, the other turns speed into skid length. Their graphs are reflections over $y=x$, and their domains and ranges are swapped.}
  \item Your answer to (d) also produces a value of $d$ for \emph{negative} speeds. Why is that half
        thrown away --- and how is this the same idea as restricting $y=x^2$ to $x\ge 0$ in class
        today?
        \ansline{Negative speed is meaningless here, so only $s\ge 0$ is kept. It is the same restriction for the same reason: the full parabola $\frac{s^2}{24}$ sends both $s$ and $-s$ to the}
        \ansline{same skid length, so it cannot be undone back into a single function. Discarding the $s<0$ half is exactly what makes $\frac{s^2}{24}$ and $\sqrt{24d}$ a genuine inverse pair.}
\end{enumerate}
\end{extensionbox}

\begin{spiralbox}
\textbf{Coming next --- Lesson 6.1: $n$th Roots \& Rational Exponents.} Today the index decided
everything --- which inputs are allowed, whether the graph stops or runs forever. Next you will learn
what an index really \emph{is}: $\sqrt[n]{a}$ written as the exponent $a^{1/n}$, so that
$a^{m/n}=\sqrt[n]{a^m}$ and every exponent rule you already know suddenly applies to radicals. That
is also where ``why does an even index refuse a negative?'' gets its full algebraic answer.
\end{spiralbox}

\begin{teachernote}
\textbf{Item 1} is the lesson's transfer task. Row~2 is the one to check first: students who write a
restriction for $\sqrt[3]{x-4}$ have learned ``radical $\Rightarrow$ restricted domain'' instead of
``\emph{even index} $\Rightarrow$ restricted domain.'' Row~4 is the sign-flip; $x\ge 2$ is the
predictable wrong answer, and it is the same slip as Exit Ticket item~3 distractor~A.

\textbf{Items 2 and 3} are deliberately a matched pair --- same questions, opposite answers on almost
every line. Item~2's absolute \emph{maximum} is the one students miss, because the parent trained them
to expect a minimum; the fix is to ask which way the arm travels from the endpoint, not what the graph
``usually'' does. Item~3's last part (symmetry) is the only question of the set that requires
algebra rather than reading, and partial credit is fine for anyone who says ``it shifted up'' without
testing a value.

\textbf{Extension} (b) is worth class time: the honest answer includes the model's limitation, not
just the arithmetic. Parts (d)--(e) are the Lesson~6.7 seed, and are the same reasoning the Activity's
Tier~E built from a falling-object model --- students who did Tier~E should recognize the structure
immediately, which is the point.
\end{teachernote}

\end{document}
